% !TEX root = ../../../main.tex
\section{Series of functions}
\label{sec:analysis-i:series-of-functions}

% BEGIN TUFTE PLACEHOLDER: supplied sample, lines 1051--1076.
\begingroup\sloppy
% Replace this entire specimen block with reviewed course notes.
Each of the document class options takes one of five justification types:
\begin{description}
  \item[\docclsopt{justified}] Fully justifies the text (sets it flush left and
    right).
  \item[\docclsopt{raggedleft}] Sets the text ragged left, regardless of which
    page it falls on.
  \item[\docclsopt{raggedright}] Sets the text ragged right, regardless of
    which page it falls on.
  \item[\doccls{raggedouter}] Sets the text ragged left if it falls on the
    left-hand (verso) page of the spread and otherwise sets it ragged right.
    This is useful in conjunction with the \docclsopt{symmetric} document class
    option.
  \item[\docclsopt{auto}] If the \docclsopt{justified} document class option
    was specified, then set the text fully justified; otherwise the text is set
    ragged right.  This is the default justification option if one is not
    explicitly specified.
\end{description}

\noindent For example, 
\begin{fullwidth}
\begin{docspec}
  \doccmdnoindex{documentclass}[symmetric,justified,marginals=raggedouter]\{tufte-book\}
\end{docspec}
\end{fullwidth}
will set the body text of the document to be fully justified and all of the
margin material (sidenotes, margin notes, captions, and citations) to be flush
against the body text with ragged outer edges.

\lipsum[2]
\par\endgroup
% END TUFTE PLACEHOLDER
